\documentclass[11pt,a4paper]{article}

\usepackage[utf8]{inputenc}
\usepackage[vietnamese]{babel}
\usepackage[top=1.5cm,bottom=1.5cm,left=1.8cm,right=1.8cm]{geometry}
\usepackage{amsmath,amssymb}
\usepackage{booktabs}
\usepackage{tabularx}
\usepackage{graphicx}
\usepackage{hyperref}
\usepackage{xcolor}
\usepackage{array}
\usepackage{caption}
\usepackage{enumitem}
\usepackage{titlesec}

\titlespacing*{\section}{0pt}{5.5pt plus 1pt minus 1pt}{2.5pt plus 1pt minus 1pt}
\titlespacing*{\subsection}{0pt}{4.5pt plus 1pt minus 1pt}{2pt plus 1pt minus 1pt}

\captionsetup{skip=2pt}
\setlist{topsep=1.5pt,itemsep=1pt,partopsep=0pt,parsep=0pt}
\setlength{\abovedisplayskip}{3pt}
\setlength{\belowdisplayskip}{3pt}
\setlength{\abovedisplayshortskip}{2pt}
\setlength{\belowdisplayshortskip}{2pt}

\hypersetup{
    colorlinks=true,
    linkcolor=blue,
    citecolor=blue,
    urlcolor=blue
}

\begin{document}

\begin{center}
{\large \textbf{ĐỀ XUẤT THIẾT KẾ KIẾN TRÚC HỆ THỐNG DỰ BÁO THỜI GIAN ĐẾN CỦA XE BUÝT (ETA) DỰA TRÊN DỮ LIỆU LỚN VÀ KIẾN TRÚC LAMBDA}}\\[3pt]
{\small Học phần: Lưu trữ và Xử lý Dữ liệu Lớn \quad $\bullet$ \quad Ngày thực hiện: \today}
\end{center}
\vspace{-4pt}

\section{Đặt Vấn Đề và Bối Cảnh Thực Tế}
Trong quản lý giao thông đô thị hiện đại, bài toán ước tính thời gian xe buýt đến trạm (Estimated Time of Arrival -- ETA) đóng vai trò quyết định đến trải nghiệm của hành khách và mức độ sử dụng giao thông công cộng. Tại các đô thị lớn như Hà Nội, thời gian di chuyển của xe buýt chịu tác động mạnh bởi mật độ dòng xe biến thiên theo khung giờ, thời tiết và đặc biệt là các điểm ùn tắc phát sinh bất ngờ.

Tuy nhiên, các dịch vụ mở hiện nay (như BusMap API) chủ yếu chỉ cung cấp thông tin tọa độ tức thời (stateless snapshots) tại thời điểm người dùng truy vấn. Việc thiếu vắng kho lưu trữ hành trình dài hạn khiến hệ thống không thể học được quy luật phân phối vận tốc lịch sử theo từng cung đường. Đồng thời, nếu chỉ dựa vào vị trí xe mà không kết hợp dữ liệu ngữ cảnh giao thông bên ngoài (như tai nạn, ùn ứ trên tuyến), kết quả dự báo ETA sẽ dễ sai lệch lớn khi có sự cố. Do đó, mục tiêu của đề tài là xây dựng một hệ thống dữ liệu lớn hoàn chỉnh theo Kiến trúc Lambda, xâu chuỗi toàn diện từ khâu thu thập dữ liệu thực tế (Crawl), lưu trữ phân tán (Storage), xử lý hai luồng (Batch \& Stream) đến giải quyết bài toán dự báo ETA.

\section{Phân Tích Đặc Trưng Dữ Liệu Lớn (Mô Hình 4V)}
Dự án được xây dựng dựa trên nguồn dữ liệu thực nghiệm đã và đang được thu thập liên tục từ mạng lưới xe buýt Hà Nội kết hợp dịch vụ dữ liệu giao thông TomTom:
\begin{itemize}[leftmargin=*,noitemsep,topsep=1.5pt]
    \item \textbf{Volume (Quy mô dữ liệu):} Hệ thống theo dõi 777 xe buýt trên 49 tuyến hành lang Hà Nội, vận hành liên tục 17 giờ mỗi ngày (05:00 -- 22:00) theo chu kỳ quét 75 giây qua GitHub Actions 24/7. Mỗi ngày phát sinh hàng trăm nghìn bản ghi định vị GPS; đồng thời thu thập định kỳ các sự kiện giao thông từ TomTom Traffic Incidents, tích lũy thành hàng triệu bản ghi trên Hugging Face và Lakehouse.
    \item \textbf{Velocity (Tốc độ phát sinh):} Dòng dữ liệu định vị được thu thập theo micro-batch 15 xe với độ trễ thấp, kết hợp nguồn trực tuyến BusMap và TimBus phục vụ theo dõi thời gian thực.
    \item \textbf{Variety (Đa dạng cấu trúc):} Hệ thống tích hợp ba nguồn dữ liệu khác nhau: (1) Chuỗi thời gian viễn thám GPS từ xe buýt (bán cấu trúc JSON); (2) Dữ liệu sự cố và ùn ứ từ TomTom Traffic Incidents API (vị trí sự cố, mức độ trễ, phạm vi ảnh hưởng); (3) Dữ liệu hạ tầng mạng lưới tuyến và trạm dừng cố định (bảng quan hệ có cấu trúc).
    \item \textbf{Value (Giá trị thực tiễn):} Chuyển đổi chuỗi tọa độ rời rạc và cảnh báo sự cố thô thành thông tin hữu ích cho hành khách: dự báo thời gian chờ thực tế tại từng trạm và cảnh báo trước tình trạng ùn ứ trên tuyến.
\end{itemize}

\section{Ý Tưởng Thiết Kế Kiến Trúc Hệ Thống (Kiến Trúc Lambda)}
Hệ thống được thiết kế theo mô hình Lambda nhằm giải quyết bài toán lớn thông qua chuỗi xử lý khép kín: \textbf{Thu thập (Crawl) $\to$ Hàng đợi (Ingestion) $\to$ Lưu trữ (Storage) $\to$ Xử lý phân tán (Processing) $\to$ Phục vụ dự báo (Serving)}.

\subsection{Khâu Thu Thập và Tiếp Nhận Dữ Liệu (Ingestion Pipeline)}
\begin{itemize}[leftmargin=*,noitemsep,topsep=1.5pt]
    \item \textbf{Thu thập dữ liệu thực tế (Crawl Layer):} Bộ thu thập chạy tiếp sức 24/7 trên GitHub Actions (phiên 9 phút) truy vấn tọa độ 777 xe buýt bằng raw SSL socket, đồng thời quét API sự cố TomTom Traffic Incidents dọc theo các hành lang tuyến.
    \item \textbf{Đệm phân vùng \& Đồng bộ Lake (Hugging Face \& Local Spool):} Dữ liệu từng phiên được đóng gói JSONL đẩy lên Hugging Face Dataset và spool cục bộ kèm kiểm tra động học (lọc \texttt{HANOI\_BBOX}, \texttt{GPS\_DRIFT\_JUMP}, \texttt{IDLE\_DEPOT}) trước khi phân phối sang Batch View và HDFS.
\end{itemize}

\newpage
\subsection{Khâu Lưu Trữ Phân Tán Bất Biến (Distributed Storage Layer)}
\begin{itemize}[leftmargin=*,noitemsep,topsep=1.5pt]
    \item \textbf{Kiến trúc Medallion Lakehouse (Bronze -- Silver -- Gold):} Dữ liệu thô từ Hugging Face được nhập bất biến vào cụm Hadoop HDFS (\texttt{/bus-eta/v1/bronze/hf\_raw}) kèm SHA-256 receipt và xuất song song ra Parquet nén Snappy phân vùng theo tuyến/ngày (\texttt{route\_code=.../dt=YYYY-MM-DD}) phục vụ Partition Pruning.
    \item \textbf{Kho vận hành SQLite (\texttt{hanoi\_bus\_system.db}):} Lưu trữ danh mục tuyến/trạm GTFS chuẩn, dữ liệu chiếu Frenet và bảng baseline phục vụ tra cứu nhanh cho tầng Web Serving.
\end{itemize}

\subsection{Khâu Xử Lý Dữ Liệu (Processing Layer: 2 Luồng Thực Thi)}
Kiến trúc phân tách rạch ròi hai luồng xử lý bổ trợ lẫn nhau:
\begin{enumerate}[leftmargin=*,noitemsep,topsep=1.5pt]
    \item \textbf{Luồng Xử lý Lô (Batch Layer -- PySpark \& PyArrow):}
    \begin{itemize}[leftmargin=*,noitemsep]
        \item Chạy trên dữ liệu lịch sử HDFS/Lakehouse: lọc nhiễu GPS, tách chuyến liên tục (\texttt{trip\_segmenter}), chiếu tọa độ 1D Frenet (\texttt{cross\_track\_m} $\le$ 35\,m) và liên kết sự cố TomTom Incidents.
        \item Xây dựng \textbf{Ma trận thời gian di chuyển cơ sở (Baseline Matrix)} theo phân vị ($p_{10}, p_{50}, p_{90}$) và huấn luyện mô hình \textbf{Residual XGBoost (\texttt{causal-eta.v1})} với 36 đặc trưng nhân quả.
    \end{itemize}
    \item \textbf{Luồng Xử lý Thời gian thực (Speed Layer -- FastAPI \& RouteHMM):}
    \begin{itemize}[leftmargin=*,noitemsep]
        \item Hợp nhất trực tuyến dữ liệu BusMap và TimBus Central API theo từng tuyến, lọc quỹ đạo nhân quả bằng mô hình ẩn Markov (\texttt{RouteHMM}).
        \item Cập nhật đồng bộ bộ trích xuất 36 đặc trưng nhân quả (\texttt{CausalFeatures}) dùng chung mã nguồn với luồng huấn luyện để tạo vector trạng thái thời gian thực.
    \end{itemize}
\end{enumerate}

\subsection{Khâu Giải Quyết Bài Toán Dự Báo ETA (Serving Layer \& Solution Idea)}
\begin{itemize}[leftmargin=*,noitemsep,topsep=1.5pt]
    \item \textbf{Mô hình Residual ETA nhân quả:} Khi người dùng tra cứu thời gian xe đến trạm đích, \texttt{ModelRuntime} tính toán phần dư điều chỉnh từ 36 đặc trưng trực tuyến cộng vào mỏ neo lịch sử:
    \begin{equation}
        \text{ETA} = \max\bigl(5.0,\ \text{baseline\_eta\_p50\_sec} + \Delta_{\text{XGBoost}}(\text{CausalFeatures}_{36})\bigr)
    \end{equation}
    \item \textbf{Chính sách ổn định \& Fallback an toàn:} Với tuyến ngoài miền huấn luyện, hệ thống tự động fallback về Baseline Matrix; tầng \texttt{eta\_policy} áp dụng các quy tắc ổn định hiển thị (pass guard, proximity cap) trước khi trả về Web Demo FastAPI.
\end{itemize}

\end{document}
